\begin{lemma}
Under the hypotheses of \(\noderef{HugeFamilySizeBound}\), assume also
\(\varepsilon_1\ge0\), and suppose the sample parameters are the rounded
hierarchy parameters
\[
m=\noderef{TwoBiteNaturalReducedVertexCount}(n),
\qquad
p=\noderef{TwoBiteEdgeProbability}\!\left(\frac12,n\right).
\]
Let
\[
H_R=H_I\cap V_R,\qquad H_B=H_I\cap V_B,
\]
and define the opposite-color masses and projection deficits by
\[
T_R=\sum_{x\in H_R}|\pi_B(X_x(I))|,\qquad
D_R=k-|\pi_R(I)|,
\]
\[
T_B=\sum_{x\in H_B}|\pi_R(X_x(I))|,\qquad
D_B=k-|\pi_B(I)|.
\]
Then every red opposite block satisfies
\[
|\pi_B(X_x(I))|\le(1+\varepsilon_2)pn\qquad (x\in H_R),
\]
and every blue opposite block satisfies
\[
|\pi_R(X_x(I))|\le(1+\varepsilon_2)pn\qquad (x\in H_B).
\]
Moreover the masses are controlled at the correct projection deficits:
\[
T_R\le(1+\varepsilon_2)D_R,\qquad
T_B\le(1+\varepsilon_2)D_B,
\]
and the corresponding one-block and two-block numerical comparisons hold:
\[
\binom{T_R}{2}_{\mathbb R}
\le
(1+\varepsilon_1)\binom{D_R}{2}_{\mathbb R},
\]
if \(T_R\le(1+\varepsilon_2)pn\), then
\[
\binom{T_R}{2}_{\mathbb R}
\le
(1+\varepsilon_1)\left(
\binom{pn}{2}_{\mathbb R}
+\binom{D_R-pn}{2}_{\mathbb R}\right),
\]
and if \((1+\varepsilon_2)pn\le T_R\), then
\[
\binom{(1+\varepsilon_2)pn}{2}_{\mathbb R}
+\binom{T_R-(1+\varepsilon_2)pn}{2}_{\mathbb R}
\le
(1+\varepsilon_1)\left(
\binom{pn}{2}_{\mathbb R}
+\binom{D_R-pn}{2}_{\mathbb R}\right),
\]
and the same conditional estimates with \(R\) and \(B\) interchanged.
\end{lemma}
\begin{proof}
First match the notation with the Lean statement.  The red and blue huge
predicates are
\[
\mathsf{hugeRed}(x)\iff
\noderef{TwoBiteHugeClass}(\omega,I,x)\wedge
  \exists r,\ x=\operatorname{inl}(r),
\]
and
\[
\mathsf{hugeBlue}(x)\iff
\noderef{TwoBiteHugeClass}(\omega,I,x)\wedge
  \exists b,\ x=\operatorname{inr}(b).
\]
Thus
\[
H_R=\operatorname{univ}.\operatorname{filter}(\mathsf{hugeRed}),
\qquad
H_B=\operatorname{univ}.\operatorname{filter}(\mathsf{hugeBlue})
\]
are exactly the finite filters in the Lean term.  The Lean quantities called
`redMass', `blueMass', `redDeficit', `blueDeficit', and `blockScale' are,
respectively,
\[
T_R,\qquad T_B,\qquad D_R,\qquad D_B,\qquad b=pn.
\]
All cardinalities in these expressions are then cast to \(\mathbb R\).

We prove the red estimates first.  Put
\[
H=H_R,\qquad h=|H|,\qquad
A_x=\pi_B(X_x(I)),\qquad
U=\bigcup_{x\in H}A_x,\qquad
W=\bigcup_{x\in H}X_x(I),
\]
and write
\[
s=\sum_{x\in H}|X_x(I)|,\qquad T=\sum_{x\in H}|A_x|,
\qquad D=D_R.
\]
If \(H=\emptyset\), then \(T=0\) and the displayed mass inequality is
immediate from \(|\pi_R(I)|\le |I|=k\) and \(0\le\varepsilon_2\).  Indeed
\[
D=k-|\pi_R(I)|\ge0,\qquad (1+\varepsilon_2)D\ge0,
\]
so \(T=0\le(1+\varepsilon_2)D\).  The pointwise red block bound is vacuous,
because no element satisfies \(\mathsf{hugeRed}\).

It remains, in this empty red case, to check the three displayed
\(\noderef{RealChooseTwo}\) comparisons.  Put \(b=pn\).  The rounded
parameter and hierarchy hypotheses give \(b\ge1\): indeed the hierarchy gives
\(1\le2pm\), the rounded-parameter comparison gives
\(0<m\le n/(\log n)^2\), and \(n>n_0\) is in the eventual range where
\((\log n)^2\ge2\), so
\[
b=pn=\frac{n}{2m}(2pm)\ge\frac{(\log n)^2}{2}\ge1.
\]
The deficit \(D=k-|\pi_R(I)|\) is a nonnegative integer.  Hence
\(\binom{0}{2}_{\mathbb R}=0\le(1+\varepsilon_1)\binom D2_{\mathbb R}\),
because \(\varepsilon_1\ge0\) and \(\binom D2_{\mathbb R}\ge0\) for an
integer \(D\ge0\).

For the small-mass two-block implication, the antecedent gives no additional
difficulty because the left side is again \(\binom02_{\mathbb R}=0\).  Since
\(\varepsilon_1\ge0\), it therefore suffices to prove
\[
\binom b2_{\mathbb R}+\binom{D-b}{2}_{\mathbb R}\ge0. \tag{E_R}
\]
If \(D=0\), then the left side of \((E_R)\) is
\[
\frac{b(b-1)}2+\frac{(-b)(-b-1)}2=b^2\ge0.
\]
If \(D=1\), it is
\[
\frac{b(b-1)}2+\frac{(1-b)(-b)}2=b(b-1)\ge0
\]
by \(b\ge1\).  If \(D\ge2\), then completing the square gives
\[
\binom b2_{\mathbb R}+\binom{D-b}{2}_{\mathbb R}
=\left(b-\frac D2\right)^2+\frac{D(D-2)}4\ge0.
\]
Thus the small-mass implication is proved in all integer cases for \(D\).

For the large-mass implication, assume its antecedent
\((1+\varepsilon_2)b\le T=0\).  Since \(0\le\varepsilon_2\) and \(b\ge1\),
we have \((1+\varepsilon_2)b\ge1\), a contradiction.  Hence the large-mass
implication is vacuous in the empty red case.
Thus assume \(H\ne\emptyset\).

First record the relative numerical comparisons supplied by
\(\noderef{ParameterHierarchyEventualComparisons}\) at the present \(n>n_0\).
The hypotheses \(m=\noderef{TwoBiteNaturalReducedVertexCount}(n)\) and
\(p=\noderef{TwoBiteEdgeProbability}(\frac12,n)\) identify the sample's
parameters with the rounded \(m\) and \(p\) that occur in that package, so the
following inequalities involve the same \(p,m\) as \(\omega\).
With
\[
t_1=\noderef{TwoBiteHugeCutoff}(n),\qquad L=\log n,
\qquad \delta=\varepsilon_2/100,
\]
the hierarchy gives \(0\le\delta\le1/100\) and the finite comparisons
\[
2pm\le\delta t_1, \tag{1}
\]
\[
\left(2k/t_1+1\right)100L^3\le\delta t_1, \tag{2}
\]
\[
\left(2k/t_1+1\right)200L^3\le\delta t_1, \tag{3}
\]
and
\[
\left(2k/t_1+1\right)400L^5\le\delta t_1. \tag{4}
\]
These are the explicit finite replacements for the paper's \(o(k)\) terms in
the huge-opposite projection argument.

The pointwise largest-block bound is direct.  If \(x\in H\), then
\[
A_x\subseteq \pi_B(N_x),
\]
because \(X_x(I)\subseteq N_x\) by \(\noderef{TwoBiteX}\).  The image
\(\pi_B(N_x)\) has cardinality at most \(|N_x|\), since a finite image cannot
have more elements than its domain, so
\[
|A_x|\le |\pi_B(N_x)|\le |N_x|.
\]
The lifted-neighborhood size clause of \(\noderef{FiberAndDegreeEvent}\) gives
\[
\bigl||N_x|-pn\bigr|\le\varepsilon_2pn,
\]
hence the upper side of the absolute-value inequality gives
\[
|N_x|\le pn+\varepsilon_2pn=(1+\varepsilon_2)pn.
\]
Therefore \(|A_x|\le(1+\varepsilon_2)pn\), which is the first red pointwise
conjunct of the Lean statement.

We next compare \(W\) with the red projection deficit.  By
\(\noderef{HugeFamilySizeBound}\), instantiated with the same
\(\omega,I,k,n,\varepsilon,\varepsilon_1,\varepsilon_2,n_0\), the same
\(\noderef{ParameterHierarchyEventualComparisons}\), the same \(n_0<n\), the
same equality \(k=\noderef{TwoBiteNaturalIndependenceScale}(1+\varepsilon,n)\),
the same \(|I|=k\), and the same
\(\noderef{FiberAndDegreeEvent}(\omega,\varepsilon_2)\), the full huge filter
has size at most \(2k/t_1\).  Since \(H\subseteq H_I\), this gives
\[
h\le 2k/t_1. \tag{5}
\]
Since every \(x\in H\) is huge, \(|X_x(I)|>t_1\), so
\[
s\ge h t_1. \tag{6}
\]
The lifted-codegree clause of \(\noderef{FiberAndDegreeEvent}\) gives
\[
|X_x(I)\cap X_y(I)|\le100L^3
\qquad(x\ne y,\ x,y\in H).
\]
We use the finite union lower bound
\[
|W|\ge
\sum_{x\in H}|X_x(I)|
  -\sum_{\{x,y\}\subset H}|X_x(I)\cap X_y(I)|,
\]
which is inclusion-exclusion truncated after pairwise intersections; no
disjointness of the \(X_x(I)\) is assumed.
Adding the sets one at a time, the loss at the step inserting \(X_x(I)\) is at
most the sum of \(|X_x(I)\cap X_y(I)|\) over the previously inserted
\(y\)'s.  There are at most \(h-1\le h\le2k/t_1\) such previous sets, and it
is harmless to bound this by \(2k/t_1+1\).  Therefore the total overlap lost in
forming \(W\) is at most
\[
\sum_{x\in H}\bigl(2k/t_1+1\bigr)100L^3
\le h\,\delta t_1
\le \delta s,
\]
where the middle inequality is exactly (2), and the last inequality is (6).
Therefore
\[
|W|\ge(1-\delta)s. \tag{7}
\]

Let \(R_W=\pi_R(W)\).  For \(x\in H\), write \(x=\operatorname{inl}(r)\).
If \(u\in\pi_R(X_x(I))\), then \(u=\pi_R(z)\) for some
\(z\in X_x(I)\).  Since \(X_x(I)\subseteq N_x\) by \(\noderef{TwoBiteX}\),
the base vertex \(u\) is red-adjacent to \(r\).  Thus
\(\pi_R(X_x(I))\) is contained in the red base-neighborhood of \(r\).
The base-degree clause of \(\noderef{FiberAndDegreeEvent}\), together with
\(0\le\varepsilon_2\le1\), bounds that neighborhood by
\((1+\varepsilon_2)pm\le2pm\).  Hence
\[
|\pi_R(X_x(I))|\le2pm. \tag{8a}
\]
Using (1) and (6),
\[
|R_W|\le\sum_{x\in H}|\pi_R(X_x(I))|
\le h\,2pm
\le \delta h t_1
\le \delta s. \tag{8}
\]
The first inequality is the finite union bound for
\(R_W=\bigcup_{x\in H}\pi_R(X_x(I))\); no disjointness is used here either.
Since \(W\subseteq I\),
\[
|\pi_R(I)|\le |\pi_R(W)|+|I\setminus W|
  = |R_W|+k-|W|.
\]
Rearranging and using (7)--(8),
\[
D=k-|\pi_R(I)|\ge |W|-|R_W|
\ge (1-2\delta)s. \tag{9}
\]
Since also \(|W|\le s\), (9) gives \(|W|\le D/(1-2\delta)\).  We also record
a lower bound that will justify the later use of monotonicity for
\(\binom{x}{2}_{\mathbb R}\).
The hierarchy comparison \(1\le2pm\), together with (1), gives
\[
1\le \delta t_1.
\]
Since \(0\le\varepsilon_2\le1\), we have \(0\le\delta\le1/100\); if
\(\delta=0\), the inequalities \(1\le2pm\le\delta t_1\) are inconsistent and
there is nothing to prove.  Otherwise \(t_1\ge1/\delta\), and because
\(H\ne\emptyset\), (6) and (9) give
\[
D\ge (1-2\delta)s>(1-2\delta)t_1
\ge (1-2\delta)/\delta\ge 98.
\]
Thus, in the non-vacuous case, \(D\ge2\).

It remains to compare the opposite projected mass \(T\) with \(W\).  The blue
fiber-size clause of \(\noderef{FiberAndDegreeEvent}\), with
\(0\le\varepsilon_2\le1\), says that every blue fiber has size at most
\(2L^2\): the clause says the fiber size is within relative error
\(\varepsilon_2\) of \(L^2\), hence is at most
\((1+\varepsilon_2)L^2\le2L^2\).  Since \(X_x(I)\) is covered by the blue
fibers indexed by \(A_x\),
\[
|X_x(I)|\le 2L^2|A_x|.
\]
Thus \(|A_x|\ge |X_x(I)|/(2L^2)\), and by (6)
\[
T=\sum_{x\in H}|A_x|\ge \frac{s}{2L^2}\ge \frac{h t_1}{2L^2}. \tag{10}
\]
The projected lifted-codegree clause of \(\noderef{FiberAndDegreeEvent}\)
gives
\[
|A_x\cap A_y|\le200L^3
\qquad(x\ne y,\ x,y\in H).
\]
Inclusion-exclusion for the projected sets gives
\[
|U|\ge T-\sum_{\{x,y\}\subset H}|A_x\cap A_y|.
\]
As above, each set \(A_x\) meets the previously inserted projected sets in at
most \((2k/t_1+1)200L^3\) points, so the overlap term is at most
\[
h(2k/t_1+1)200L^3
\]
Rewrite this upper bound as
\[
\frac{h}{2L^2}\,(2k/t_1+1)400L^5,
\]
and using the hierarchy comparison (4), this is at most
\[
\delta\frac{h t_1}{2L^2}.
\]
The lower bound (10) says \(T\ge h t_1/(2L^2)\), so the overlap term is at
most \(\delta T\).  Hence
\[
|U|\ge(1-\delta)T.
\]
The set \(U\) is the blue projection of \(W\), so \(|U|\le |W|\).  Combining
this with (9), and using \(1-\delta>0\) and \(1-2\delta>0\), gives
\[
T\le\frac{|U|}{1-\delta}
\le\frac{|W|}{1-\delta}
\le\frac{D}{(1-\delta)(1-2\delta)}
\le(1+\varepsilon_2)D. \tag{11}
\]
The last inequality uses \(\delta=\varepsilon_2/100\) and
\(0\le\varepsilon_2\le1\); explicitly,
\[
\frac{1}{(1-\delta)(1-2\delta)}\le1+\varepsilon_2
\]
for \(0\le\delta=\varepsilon_2/100\le1/100\).

We now invoke \(\noderef{HugeOppositeProjectionMassNumerics}\) for the real
binomial comparisons.  Set
\[
b=pn,\qquad M=(1+\varepsilon_2)b.
\]
The side conditions in that helper have all been established here.  The
eventual-comparisons package gives
\[
0\le\varepsilon_2,\qquad \varepsilon_2\le1,\qquad
3\varepsilon_2\le\varepsilon_1/10,
\]
and \(\varepsilon_1\ge0\) is one of the hypotheses of the present lemma.  The
definition of \(p=\noderef{TwoBiteEdgeProbability}(\frac12,n)\) gives
\(b=pn\ge0\).  The deficit
\[
D=k-|\pi_R(I)|
\]
is a nonnegative integer, because the finite image \(\pi_R(I)\) has
cardinality at most \(|I|=k\), so \(D\) is the cast of the natural difference
\(k-|\pi_R(I)|\).  The mass
\[
T=\sum_{x\in H_R}|\pi_B(X_x(I))|
\]
is also a nonnegative integer, being a finite sum of cardinalities.  In the
nonempty case above, (9) and the subsequent lower-bound argument give
\(D\ge98\), and (11) gives
\[
0\le T\le(1+\varepsilon_2)D.
\]
This numerical helper is invoked only in the present nonempty case; the
empty red case was discharged directly above precisely because the helper's
large-deficit hypothesis need not hold when \(H_R=\emptyset\).  Therefore
\(\noderef{HugeOppositeProjectionMassNumerics}\), applied with
\((D,b,T)=(D_R,pn,T_R)\), gives exactly the three red numerical claims:
\[
\binom{T_R}{2}_{\mathbb R}
\le
(1+\varepsilon_1)\binom{D_R}{2}_{\mathbb R},
\]
\[
T_R\le(1+\varepsilon_2)pn\Longrightarrow
\binom{T_R}{2}_{\mathbb R}
\le
(1+\varepsilon_1)\left(
\binom{pn}{2}_{\mathbb R}+\binom{D_R-pn}{2}_{\mathbb R}\right),
\]
and
\[
(1+\varepsilon_2)pn\le T_R\Longrightarrow
\binom{(1+\varepsilon_2)pn}{2}_{\mathbb R}
+\binom{T_R-(1+\varepsilon_2)pn}{2}_{\mathbb R}
\le
(1+\varepsilon_1)\left(
\binom{pn}{2}_{\mathbb R}+\binom{D_R-pn}{2}_{\mathbb R}\right).
\]
The helper is where the \(0/1\) regimes for
\(\binom{x}{2}_{\mathbb R}\) are handled: monotonicity is used only on
\([1,\infty)\), and the exceptional interval \(0<D_R-pn<1\) is controlled
using the fact that both \(D_R\) and \(T_R\) are integer-valued.

We now prove the blue estimates with the projection choices written out.  Put
\[
H'=H_B,\qquad h'=|H'|,\qquad
B_x=\pi_R(X_x(I)),\qquad
U'=\bigcup_{x\in H'}B_x,\qquad
W'=\bigcup_{x\in H'}X_x(I),
\]
and write
\[
s'=\sum_{x\in H'}|X_x(I)|,\qquad
T'=T_B=\sum_{x\in H'}|B_x|,\qquad
D'=D_B=k-|\pi_B(I)|.
\]

If \(H'=\emptyset\), then \(T'=0\).  Since
\(|\pi_B(I)|\le |I|=k\), the blue deficit \(D'=D_B\) is nonnegative, and
\[
T'=0\le(1+\varepsilon_2)D'
\]
because \(0\le\varepsilon_2\).  The pointwise blue block bound is vacuous,
because no element satisfies \(\mathsf{hugeBlue}\).  The same hierarchy
calculation as above gives
\[
b=pn=\frac{n}{2m}(2pm)\ge\frac{(\log n)^2}{2}\ge1.
\]
The deficit \(D_B=k-|\pi_B(I)|\) is a nonnegative integer, so
\[
\binom{0}{2}_{\mathbb R}
\le(1+\varepsilon_1)\binom{D_B}{2}_{\mathbb R}
\]
because \(\varepsilon_1\ge0\) and
\(\binom{D_B}{2}_{\mathbb R}\ge0\) for an integer \(D_B\ge0\).
For the small-mass implication, the left side is again
\(\binom02_{\mathbb R}=0\).  It remains only to prove
\[
\binom b2_{\mathbb R}+\binom{D_B-b}{2}_{\mathbb R}\ge0. \tag{E_B}
\]
If \(D_B=0\), then the left side of \((E_B)\) is
\[
\frac{b(b-1)}2+\frac{(-b)(-b-1)}2=b^2\ge0.
\]
If \(D_B=1\), it is
\[
\frac{b(b-1)}2+\frac{(1-b)(-b)}2=b(b-1)\ge0
\]
by \(b\ge1\).  If \(D_B\ge2\), then
\[
\binom b2_{\mathbb R}+\binom{D_B-b}{2}_{\mathbb R}
=\left(b-\frac{D_B}{2}\right)^2+\frac{D_B(D_B-2)}4\ge0.
\]
Finally, the large-mass antecedent
\((1+\varepsilon_2)b\le T'=0\) contradicts
\((1+\varepsilon_2)b\ge b\ge1\).  Thus the large-mass implication is
vacuous.  This proves all blue conjuncts when \(H_B=\emptyset\).

Assume now that \(H_B\ne\emptyset\).  The pointwise cap uses the opposite
projection for a blue huge vertex.  If \(x\in H_B\), then
\[
B_x=\pi_R(X_x(I))\subseteq \pi_R(N_x),
\]
because \(X_x(I)\subseteq N_x\).  A finite image has no more points than its
domain, so \(|B_x|\le |\pi_R(N_x)|\le |N_x|\).  The lifted-neighborhood size
clause of \(\noderef{FiberAndDegreeEvent}\) gives
\[
\bigl||N_x|-pn\bigr|\le\varepsilon_2pn,
\]
and hence
\[
|B_x|\le |N_x|\le(1+\varepsilon_2)pn.
\]
This is exactly the pointwise blue conjunct in the Lean statement.

The full huge filter bound from \(\noderef{HugeFamilySizeBound}\), applied
with the same data as before, gives
\[
h'\le2k/t_1, \tag{12}
\]
because \(H_B\subseteq H_I\).  Since each element of \(H_B\) is huge,
\[
s'\ge h't_1. \tag{13}
\]
The blue-blue lifted-codegree clause of \(\noderef{FiberAndDegreeEvent}\)
gives
\[
|X_x(I)\cap X_y(I)|\le100L^3
\qquad(x\ne y,\ x,y\in H_B).
\]
Using the same finite inclusion-exclusion lower bound as in the red case,
\[
|W'|\ge
\sum_{x\in H'}|X_x(I)|
  -\sum_{\{x,y\}\subset H'}|X_x(I)\cap X_y(I)|,
\]
and then using (2), (12), and (13), the total overlap loss is at most
\(\delta s'\).  Hence
\[
|W'|\ge(1-\delta)s'. \tag{14}
\]

Let \(B_W=\pi_B(W')\).  For \(x\in H_B\), the same-color projection
\(\pi_B(X_x(I))\) is handled symmetrically.  Write
\(x=\operatorname{inr}(b_0)\).  If \(v\in\pi_B(X_x(I))\), then
\(v=\pi_B(z)\) for some \(z\in X_x(I)\), and the containment
\(X_x(I)\subseteq N_x\) shows that \(v\) lies in the blue base-neighborhood
of \(b_0\).  The blue base-degree clause of
\(\noderef{FiberAndDegreeEvent}\), together with
\(0\le\varepsilon_2\le1\), bounds this neighborhood by
\((1+\varepsilon_2)pm\le2pm\).  Hence
\[
|\pi_B(X_x(I))|\le2pm. \tag{15a}
\]
Therefore, by the finite union bound and the hierarchy comparison (1),
\[
|B_W|
\le\sum_{x\in H'}|\pi_B(X_x(I))|
\le h'\,2pm
\le \delta h't_1
\le \delta s'. \tag{15}
\]
Since \(W'\subseteq I\),
\[
|\pi_B(I)|\le |\pi_B(W')|+|I\setminus W'|
  =|B_W|+k-|W'|.
\]
Rearranging and using (14)--(15) gives
\[
D_B=k-|\pi_B(I)|
\ge |W'|-|B_W|
\ge(1-2\delta)s'. \tag{16}
\]
Also \(|W'|\le s'\), so \(|W'|\le D_B/(1-2\delta)\).  As in the red case,
the inequalities \(1\le2pm\le\delta t_1\), \(H_B\ne\emptyset\), (13), and
(16) give
\[
D_B\ge(1-2\delta)s'>(1-2\delta)t_1
\ge(1-2\delta)/\delta\ge98. \tag{17}
\]

It remains to control the blue opposite mass \(T_B\).  Now the relevant
fibers are red fibers, because \(B_x=\pi_R(X_x(I))\).  The red fiber-size
clause of \(\noderef{FiberAndDegreeEvent}\), with
\(0\le\varepsilon_2\le1\), gives every red fiber size at most \(2L^2\).
Since \(X_x(I)\) is covered by the red fibers indexed by \(B_x\),
\[
|X_x(I)|\le2L^2|B_x|.
\]
Thus
\[
T_B=\sum_{x\in H'}|B_x|
\ge\frac{s'}{2L^2}\ge\frac{h't_1}{2L^2}. \tag{18}
\]
For the projected union \(U'=\bigcup_{x\in H'}B_x\), the projected
lifted-codegree clause used is the blue-blue clause with red projections:
\[
|B_x\cap B_y|\le200L^3
\qquad(x\ne y,\ x,y\in H_B).
\]
Finite inclusion-exclusion gives
\[
|U'|\ge T_B-\sum_{\{x,y\}\subset H'}|B_x\cap B_y|.
\]
Each \(B_x\) meets the previously inserted projected sets in at most
\((2k/t_1+1)200L^3\) points, so the overlap term is at most
\[
h'(2k/t_1+1)200L^3
=\frac{h'}{2L^2}(2k/t_1+1)400L^5.
\]
By (4), this is at most
\[
\delta\frac{h't_1}{2L^2},
\]
and by (18) it is at most \(\delta T_B\).  Hence
\[
|U'|\ge(1-\delta)T_B. \tag{19}
\]
The set \(U'\) is the red projection of \(W'\), so \(|U'|\le |W'|\).  Combining
(16), \(|W'|\le s'\), and (19), and using the same elementary inequality
\[
\frac{1}{(1-\delta)(1-2\delta)}\le1+\varepsilon_2,
\]
gives
\[
T_B\le(1+\varepsilon_2)D_B. \tag{20}
\]

We apply \(\noderef{HugeOppositeProjectionMassNumerics}\) with
\[
(D,b,T)=(D_B,pn,T_B).
\]
The side conditions are now explicit: \(\varepsilon_1\ge0\) is a hypothesis;
\(0\le\varepsilon_2\le1\) and
\(3\varepsilon_2\le\varepsilon_1/10\) come from
\(\noderef{ParameterHierarchyEventualComparisons}\); \(b=pn\ge0\) follows
from the rounded definition of \(p\); \(T_B\ge0\) is a finite sum of
cardinalities; \(D_B=k-|\pi_B(I)|\) is a nonnegative integer because
the finite image \(\pi_B(I)\) has cardinality at most \(|I|=k\), so \(D_B\)
is the cast of the natural difference \(k-|\pi_B(I)|\); \(T_B\) is
integer-valued as a finite sum of cardinalities; (17) gives \(98\le D_B\);
and (20) gives
\[
0\le T_B\le(1+\varepsilon_2)D_B.
\]
As in the red case, the helper is used only after the nonempty blue filter
has supplied the large-deficit condition; the empty blue filter was already
settled directly.
The numerical helper therefore gives the three blue binomial conclusions
\[
\binom{T_B}{2}_{\mathbb R}
\le
(1+\varepsilon_1)\binom{D_B}{2}_{\mathbb R},
\]
\[
T_B\le(1+\varepsilon_2)pn\Longrightarrow
\binom{T_B}{2}_{\mathbb R}
\le
(1+\varepsilon_1)\left(
\binom{pn}{2}_{\mathbb R}+\binom{D_B-pn}{2}_{\mathbb R}\right),
\]
and
\[
(1+\varepsilon_2)pn\le T_B\Longrightarrow
\binom{(1+\varepsilon_2)pn}{2}_{\mathbb R}
+\binom{T_B-(1+\varepsilon_2)pn}{2}_{\mathbb R}
\le
(1+\varepsilon_1)\left(
\binom{pn}{2}_{\mathbb R}+\binom{D_B-pn}{2}_{\mathbb R}\right).
\]
Together with the pointwise cap and mass bound already proved, these are
exactly the five blue conjuncts of the Lean statement.  The five red conjuncts
were proved above, so the lemma follows.
\end{proof}
